\noindent\textbf{Elección del procedimiento.} La masa que atraviesa el disco por segundo es $\rho AV$. Su energía cinética específica es $V^2/2$, por lo que $P_{viento}=\rho AV^3/2$. El modelo ideal de disco actuador limita la fracción extraíble a $16/27$, dato teórico usado en el capítulo.

\begin{enumerate}

\item Convertir diámetro del rotor en área de barrido.

\[\fitmath{A=\frac{\pi\cdot {D}^{2}}{4}}\]

\[\fitmath{A=\frac{\pi\cdot {(120)}^{2}}{4}\approx 11309{,}73\,\mathrm{m^{2}}}\]

\item Multiplicar caudal másico por energía cinética específica.

\[\fitmath{\mathrm{Pwind}=\frac{\rho\cdot A\cdot {V}^{3}}{2}}\]

\[\fitmath{\mathrm{Pwind}=\frac{(1{,}225)\cdot (11309{,}73355)\cdot {(10)}^{3}}{2}\approx 6927212\,\mathrm{W}}\]

\item Aplicar el límite ideal, no un rendimiento real medido.

\[\fitmath{\mathrm{Pmax}=\frac{16\cdot \mathrm{Pwind}}{27}}\]

\[\fitmath{\mathrm{Pmax}=\frac{16\cdot (6927211{,}801)}{27}\approx 4105014\,\mathrm{W}}\]

\end{enumerate}
